\section{Methods}\label{sec:methods}

\subsection{Pre-Specification and the Development Round}\label{sec:prespec}

The hypotheses, the design of the runs, the statistics and the rules of decision were pre-specified
in the paper's repository at commit \PreSpecCommit, before the run. From then on the file was frozen,
and every change went into its revision log with its reason. Section~\ref{sec:deviation} describes the
log's entries. The analysis that decides every hypothesis was written and
tested on synthetic runs before the freeze. The two files of macros that carry its numbers into this
paper were extended after the run, from its outputs and raw data, without changing a decision.

A development round measured a pre-release of the same lookup code, with the same harness and
design, on other seeds. It guided the changes that Section~\ref{sec:rework} describes. Its numbers
are not reported, and none of its seeds was reused. The seeds of the run had timed nothing before
it. They had served only to check answers, in the agreement tests and the sanitizer records.

\subsection{Host and Builds}\label{sec:host}

Every measurement ran on one lab host, L: an \HostLCpu{} (\HostLArch) with Linux \KernelVersion. The
pre-specification froze the host, with no package update and no restart from the first sanitizer
record to the end of the last run. Every process ran on CPU~\TimedCpu, whose sibling hardware thread, CPU~\TimedCpuSibling,
ran nothing of ours. The run recorded the performance governor, boost off, a mean clock of
\MeanMhz~MHz, transparent huge pages always on, and the NMI watchdog on.

C and C++ were compiled by clang \ClangVersion, Rust by rustc \RustcVersion, and Go by Go \GoVersion.
Every arm was compiled for L's own instruction set. The C and C++ arms used \FlagsCxx, the Rust arms
\FlagRust{} and the Go arms \FlagGo. On L these resolve to the target \TargetCpu. Two competitors
point to a native build. glaze's performance guide, written for its serialization, recommends a
native build when a program runs where it was built, and uWebSockets' own example build uses one.
Fairness needs the same level for every arm, v2 included. It also keeps H2's subtraction matched,
because the harness's loop is compiled into each arm's translation unit. Each build writes the
resolved target into a header that every target compiles. A build for another instruction set
therefore has other inputs, and the gate of Section~\ref{sec:gate} refuses it.

Link-time optimization is off for every C and C++ arm, v2 included. The Rust arms use the harness's
release profile, \RustProfile. It optimizes across all crates, so it is at least cargo's default
profile and favours the Rust arms. The Go arms are built as a C archive without
profile-guided optimization, which would need a profile of a workload that this design does not
have. Appendix~\ref{app:config} gives each competitor's configuration and the documentation it
follows.

\subsection{Arms}\label{sec:arms}

An arm is one router behind the harness's common interface. All arms are linked into one binary.

\runin{v2.} \arm{regexmatcher-v2} adds the routes at run time, builds the table, and looks up
through one out-of-line function, \FindVtwo. \arm{regexmatcher-v2-ct} declares the same
routes as a constant table in a generated translation unit, one per shape and size, aligned to a
page. It shares its adapter with \arm{regexmatcher-v2}, takes part in H7 only, and by design refuses
tables above \SizeThousand{} routes.

\runin{Competitors.} The sixteen routers of Table~\ref{tab:survey}, each pinned by commit and by
the hash of its source archive. Each arm's timed lookup starts from the method and the path, in the
form that its library's request carries them. It ends with the route and the captured values, as the
router's interface returns them. Steps of a framework outside route lookup, such as pooling a request and writing a
response, are left out for every arm. Appendix~\ref{app:lookup} gives each arm's timed call and the
differences that remain, because removing them would leave the router's interface.

\runin{Reference arms.} v1 behind its wrapper (\arm{regexmatcher-v1}), a radix tree of the
harness, an exact-match hash table of the harness, and a sequence of \texttt{std::regex}. The hash
table and the radix tree run in every cell, with as many processes as v2. The other two run up to
$m=\RefMaxSize$, in \RefReps{} processes per cell. On the \shape{static} shape, whose one method has
only literal routes, v2 answers from its exact-match table, the mechanism of the hash arm.

\runin{Null arms.} A null arm measures the harness's loop around a lookup that matches nothing:
one out-of-line call that reads the path's length. H2 subtracts from each arm the null of its own
loop. \arm{null} answers 405 itself, so the harness adds nothing to its answer. It is the null of
v2, of the compile-time table and of the ablation arm. A second null arm, \NullNoFive, has no 405 of
its own, so the harness applies its 405 rule to it, as to every arm without one. It is the null of
the C++ competitors, of the Rust competitors behind their C functions and of the reference arms.
Each Go arm runs a null pass of its own, in its own process, over the same data.

\runin{The ablation arm.} \AblationArm{} is the pre-release's header with its namespace and macro
prefix renamed and nothing else changed. It runs in every cell of the main grid, decides nothing, and
shows what the rework of Section~\ref{sec:rework} changed.

\subsection{Tables and Rings}\label{sec:tables}

A cell is one arm, one table shape and one table size $m$. Table~\ref{tab:shapes} lists the shapes.
The sizes are $m=\SizeTen$, \SizeHundred, \SizeThousand, \SizeTenThousand{} and \SizeHundredThousand.
Literals are short pseudo-words. A route's first literal comes from a small vocabulary and every later
one from a large one. A table therefore shares a few top-level words and branches widely below them,
as many REST interfaces do. Each process draws its table from its table seed.

\begin{table}[H]
\caption{The seven table shapes. \texttt{l} is a literal and \texttt{p} a parameter; \texttt{*p} is a
catch-all.\label{tab:shapes}}
\centering\small
\begin{tabularx}{\textwidth}{@{}l l X@{}}
\toprule
\textbf{Shape} & \textbf{Segments} & \textbf{What it tests} \\
\midrule
\shape{static} & \texttt{/l/l/l} & literal routes only, one method \\
\shape{param-last} & \texttt{/l/l/\{p\}} & a parameter as the last segment \\
\shape{param-first} & \texttt{/\{p\}/l/l} & a parameter as the first segment \\
\shape{rest} & \texttt{/l/\{p\}/l/\{q\}} & two parameters between literals \\
\shape{wild} & \texttt{/l/l/\{*p\}} & a catch-all \\
\shape{mixed-disjoint} & groups of four & a fully literal route and three routes of one resource, \texttt{/b/r/\{p\}}, \texttt{/b/r/\{p\}/s} and \texttt{/b/r/\{p\}/s/\{q\}}, under different first words \\
\shape{mixed-overlap} & groups of four & \texttt{/t/r/s}, \texttt{/t/r/\{p\}}, \texttt{/t/r/s/\{p\}} and \texttt{/t/\{p\}/u/\{q\}}: precedence and backtracking \\
\bottomrule
\end{tabularx}
\end{table}

A ring holds \RingSize{} lookups drawn uniformly over the table's routes, from the process's ring
seed. Every lookup in the main grid hits a route. In \shape{mixed-overlap}, half of the queries for
the last route of a group carry the literal \texttt{r} as the parameter's value. Such a query first
descends the literal, fails below it, and must return to the parameter. We call these queries decoys.

\subsection{Measurements}\label{sec:measurements}

Each process builds one table, then measures one arm (\texttt{bench/core/runner.hpp}). It reports the
median time per lookup over nanobench's epochs. Per lookup it also reports instructions, branch
misses, cycles, L1d misses, L2 misses of data requests and L1 DTLB misses, all in user space. The
dispatch-stall cycles of every unit-mask bit of AMD's events \EventsAmd{} are counted, each bit on its
own. We found no AMD document that defines these events for L's processor, so they rest on a
validation run on L, kept as \StallEvidence. Latency percentiles come from one lookup at a time. Heap allocations per lookup are counted in a
separate pass, never in a timed loop. The process also reports the build time, the bytes allocated by
the build and the heap bytes that the built table holds.

The timed loops read a compact copy of the ring, made before the table is built. A Go arm loads its
own copy into Go's heap at its first timed pass. Go's collector and system monitor share the
measuring core with the timed loop, and the thread's counters do not count their time. The Go arms'
clock, their spread over epochs and their collections during the epochs are therefore in Table~S1.

\subsection{Checking Every Answer}\label{sec:checks}

A time counts only if every answer behind it is right. In each process, an agreement pass checks the
route and every captured value of each query against the ring's expectation. The pass stops after
\VerifyBudgetSeconds~s, once it has checked at least \VerifyMinQueries{} queries. After the last
measured cell, the run checks every process that the pass cut on every query of the same table and
ring. That check is untimed, and its tool is linked from the run's own build. It counts only if its
table and ring digests equal the process's. A process that neither check verified on every query
fails its cell.

Every process's answers are also judged against declared path semantics. A router whose answers on a
shape differ by its own semantics is measured all the same and shown as incomparable, with the range
of differing answers. It then takes no part in that shape's cells. Each declaration names the one
kind of wrong answer it allows, and any other wrong answer fails the cell. There are four:
\begin{itemize}
\item Boost.URL's example router on \shape{mixed-overlap} finds the right route but keeps a value
captured in a branch it left.
\item gin on \shape{mixed-overlap} answers 404 for each query that needs a return to an earlier
parameter. Wherever its route is right, its values are right.
\item r3 on \shape{mixed-overlap} answers 404 for such queries, and for more queries as the table
grows.
\item uWebSockets on \shape{wild} finds the route, but its catch-all returns no value. So uWebSockets
takes part in none of the \shape{wild} cells.
\end{itemize}
An arm that rejects a table while it is built is shown as refused, with the router's reason.

Eight probes test the rules of Section~\ref{sec:semantics} on every arm. The most specific route
must win in either registration order (two probes). Every pchar must be accepted in a parameter and
returned as sent. Bytes must be compared as sent, without decoding. A catch-all must take a non-empty
remainder, and an empty segment must not match a parameter. A trailing slash must be strict, and a
path that matches only under another method must give 405.

\runin{The cut for slow arms.} A timed pass that would take longer than \PassBudgetSeconds~s is
cut to the prefix of the ring that fits, never to fewer than \PassMinQueries{} queries. The cut is
decided from the agreement pass's time. The ring is in random order, but a cut arm repeats a short
prefix, which warms its predictors and caches. A cut arm's time and counters are therefore optimistic
for it. The analysis flags every cell where v2 or the fastest competitor was cut. H2 reports the
median competitor's branch misses with and without the cut competitors.

\runin{Stopping.} A process has a budget of \ProcessBudgetSeconds~s. When a cell's first process
exceeds it, the cell's other processes do not run, and the cell is reported as over the budget.

\subsection{The Sanitizer Gate}\label{sec:gate}

A number is citable only if every measured build compiled first-party inputs that green sanitizer
records cover. On L the records are AddressSanitizer with UndefinedBehaviorSanitizer,
ThreadSanitizer and MemorySanitizer, all with clang \ClangVersion. A record covers a build when the
hashes of their inputs match, with the same compiler, configuration, pins and fetched archives.
RegexMatcher's own suite, with its negative-compilation cases, passed under all three on L. That
suite's records cover v2's headers. The harness's records cover every arm, and the server's records
cover the server with its end-to-end exercise. Every measured build passed the gate, and a build of
every arm for another instruction set failed it.

MemorySanitizer needs every library in the process instrumented, and the Go and Rust standard
libraries are prebuilt and uninstrumented. Its record therefore leaves out the \MsanGapArms{} arms
behind a foreign-function boundary: matchit, actix-router, path-tree, httprouter, gin, net/http and
chi. The gap is declared in the repository. Under AddressSanitizer their Go and Rust code is
instrumented too, though not the Rust standard library. Under ThreadSanitizer only their C++
adapters are instrumented. No binary of the records links or calls io\_uring, so MemorySanitizer's blind
spot for the kernel's writes through it does not apply.

\subsection{Seeds and Order}\label{sec:seeds}

Each cell runs in $R=\RoundTwoReps$ processes. Process $k$ uses the $k$th table seed and the $k$th ring seed.
The table seeds run from \TableSeedFirst{} to \TableSeedLast{}, and the ring seeds from
\RingSeedFirst{} to \RingSeedLast. No two pairs share a table or a ring, so the pairs are
independent, and the analysis treats them so. It refuses a cell whose v2 processes do not use exactly
these pairs. All arms ran in one binary in the same run, and the order of the processes was shuffled
with the seed \RoundTwoOrderSeed.

Why $R=\RoundTwoReps$: the project's rule on replicates asks for enough of them that an exact paired
test can pass the family's Holm threshold. With sixteen independent pairs, the smallest one-sided p of the
exact sign test is one half to the power $R$, that is \SignMinP. Holm's first threshold for the family of \HOneCells{}
cells is \Alpha{} divided by \HOneCells, that is \HolmFirst. One pair at or above the margin gives
\SignOneAbove. The smallest p of the bootstrap, about \BootMinP, is also below Holm's first threshold.
Each table seed is in one pair, so the bootstrap clustered by table seed is the bootstrap over pairs.

\subsection{Statistics}\label{sec:stats}

A comparison of two arms in a cell is paired by seed pair. Its statistic is the median over the
\RoundTwoReps{} pairs of the ratio of the two arms' times. For H1, the fastest competitor of a cell is
the included competitor with the lowest median time. It may differ from cell to cell, and the
bootstrap chooses it again in every resample. Two computations, both pre-specified, test each cell.

\runin{The BCa interval.} \RoundTwoResamples{} resamples of the pairs, drawn with replacement,
keep each pair's values together. The bias correction counts ties as half, and the acceleration comes
from the jackknife over pairs. With sixteen pairs, the leave-one-out medians take two values, eight
times each. The acceleration is then zero whenever the fastest competitor is the same in every
leave-one-out set. In this run it was zero up to rounding, at most \HOneMaxAccel, so every interval is the
bias-corrected percentile interval. The one-sided p of the null hypothesis is the level $\alpha$ at
which the interval's upper bound equals the margin. That interval is two-sided, with coverage one
minus twice $\alpha$.

\runin{The exact sign test.} $X$ counts the pairs whose ratio is below the margin. A tie counts as
not below. Under the null hypothesis that the median ratio is at least the margin, $X$ is binomial
with $n=\RoundTwoReps$ and probability one half at the boundary, and $p=P(X\geq x)$ exactly.

The implementation uses Python's standard normal distribution. On L it reproduced SciPy's BCa interval
to the last digit on the same draws, and SciPy's exact binomial test. One generator per hypothesis,
seeded with \RoundTwoResamplingSeed, serves the tested cells in the family's order. An untested cell
draws nothing.

\subsection{The Hypotheses}\label{sec:hypotheses}

\runin{H1, lookup time.} In every cell in which a competitor holds the table, v2 is not slower
than the fastest of them. The null hypothesis of a cell is that the ratio is at least \Margin. Holm's
step-down runs over the \HOneCells{} cells at a family-wise \Alpha, one-sided, for each computation. A
cell passes if it passes under both. A cell that cannot be tested enters Holm with a p of one and does
not pass. H1 holds if at least \HOneNeed{} cells pass and the upper bound of the unadjusted \Level\% BCa
interval is at most \HOneCapTwo{} in every tested cell. Superiority, with the margin \SupMargin, is
tested the same way in a Holm family of its own. It decides nothing about H1. The paper calls v2
faster only in the cells that pass superiority.

\runin{H2, work per lookup.} In every cell, v2 executes no more instructions per lookup than the
fastest competitor, and takes no more branch misses than the median competitor. Instructions are
compared net of the null of each arm's own loop, and branch misses raw. Values are medians over the
processes, and the comparison is strict. A cell that cannot be judged does not hold.

\runin{H3, no allocation.} v2 makes no heap allocation per lookup in any process of the
\HOneCells{} cells, captured values included.

\runin{H4, build and memory.} In every cell, v2's build time is at most the median competitor's,
and at $m=\SizeHundredThousand$ it is below \BuildLimitSeconds~s. The heap bytes that its table holds
are at most \HeapFactor{} times the smallest competitor's. All three must hold in every cell.

\runin{H5, checked when compiled, correct when run.} (a) Each of the \NegCases{}
negative-compilation cases fails to compile, with a diagnostic that names its reason. (b) v2 passes
every probe. (c) Over HTTP/1.1, in the paper's minimal server, a HEAD request on a path that only GET
routes match is served by the GET route without content. A 405's Allow field lists HEAD whenever it
lists GET~\cite{fielding2022rfc9110}. HTTP/2 is out of scope, because the server has none.

\runin{H6, the client sees the difference.} In each of nine cells, the server answers more
requests per second with v2 than with v1. That is, the lower bound of the ratio's \Level\% BCa
interval is above \SupMargin. In each, the measured difference in time per request is at least \CarryBound{} of
the difference that the lookup times predict (Section~\ref{sec:server}).

\runin{H7, a table built at compile time is not slower.} On the same routes and ring, the
compile-time table is not slower than v2's table built at run time. In every cell, both computations
must pass at the margin \Margin, each at a one-sided \Alpha{} per cell. The sign test then needs
\HSevenSignNeed{} of \RoundTwoReps{} pairs below the margin. H7 covers the seven shapes at $m\leq
\SizeThousand$, \HSevenCells{} cells.

\runin{What would count against the claim.} The claim is that v2 is the fastest route matcher
measured here, in the setting of Section~\ref{sec:host}. ``Faster'' rests on H1's superiority family
alone, cell by cell, and only against the competitors included in a cell. H1 or H2 failing in a whole
shape narrows the claim to the other shapes, and a loss in a single cell is reported in that cell.
If H4, H6 or H7 fails, the paper reports the loss and does not make that hypothesis's claim.

\subsection{The Minimal Server and the End-to-End Run}\label{sec:server}

H5(c) and H6 use a minimal HTTP/1.1 server written for this paper, at commit \ServerCommit. It runs
one worker thread over epoll with keep-alive, and answers pipelined requests in order. Each route
answers with a fixed body of \HSixBodyBytes{} bytes. HEAD and the Allow field live above the router,
so every router arm runs the same code for them. One binary holds three router arms behind one
function pointer: v1 behind its wrapper, v2, and a null router that answers without reading the path.
The server is not a full framework. It parses requests, routes them and writes fixed responses,
and nothing else.

H6 has nine cells: the seven shapes at $m=\SizeTen$, and \shape{rest} and \shape{param-last} at
$m=\SizeTenThousand$. Each serves the routes of table seed \HSixTableSeed. Its targets are one per
route, with each parameter filled as ring seed \HSixRingSeed{} fills it, in shuffled order. The server
runs on CPU~\HSixServerCpu, whose sibling CPU~\HSixServerSibling{} idles. The load generator runs on
CPUs \HSixGenCpuFirst{} to \HSixGenCpuLast, over \HSixConnections{} connections over loopback, with
one request in flight on each. Each window starts a fresh server, probes it, warms it up for
\HSixWarmupSeconds~s and measures for \HSixWindowSeconds~s. Requests per second are the completed
successful responses over the measured wall time.

The arms run in mirrored rounds of v1, v2, null, null, v2 and v1, and one round is one pair. A window
is invalid if its server did not start or answer its probe, or if it completed no request. It is also
invalid if its error share exceeded \HSixErrorSharePct\%, if the generator's CPUs were more than
\HSixGenBusyPct\% busy, or if the clock drifted more than \HSixMhzDriftPct\%. A round with an invalid
window is not a valid pair. A cell runs at least \HSixMinPairs{} and at most \HSixMaxPairs{} rounds. It
stops once a bootstrap interval of the geometric mean ratio, for v2 and for the null arm against v1,
lies outside $[\HSixStopLow, \HSixStopHigh]$. It also stops once that interval's half-width is below
\HSixStopHalfWidthPct\%.

Per pair, each arm's rate $X$ is the mean over its two windows, and its time per request $\tau$ is
the reciprocal of $X$. The ratio is $X_{v2}/X_{v1}$, and the measured difference is
$D_{\mathrm{meas}}=\tau_{v1}-\tau_{v2}$. The lookup difference $D_{T0}=t_{v1}-t_{v2}$ is v1's time
minus v2's in the main grid's process at pair (\HSixTableSeed, \HSixRingSeed), in the same cell.
H6(a) needs the lower bound of the ratio's \Level\% BCa interval above \SupMargin{} in all nine cells.
H6(b) needs the lower bound of the interval of $D_{\mathrm{meas}}/D_{T0}$ at least \CarryBound{} in
all nine. $D_{T0}$ is one value per cell, so the interval of (b) covers the variation of the
end-to-end run only. The bound \CarryBound{} was chosen before the run, informed by development data
from an earlier build that the paper does not report. With at least \HSixMinPairs{}
pairs, a cell's exact sign test can reach \HSixSignMinP. Beside each cell, and untested, we report
the null arm's time per request. We also report the ratio that the lookup times predict,
$P=(\tau_{\mathrm{null}}+t_{v1})/(\tau_{\mathrm{null}}+t_{v2})$.

\subsection{Exploratory Runs}\label{sec:explore}

Each exploratory run used the main grid's gated build and order seed, on the first $R$ pairs. None
has an interval, and none decides anything.
\begin{itemize}
\item Miss-heavy rings, where half the queries match no route: every cell, $R=\ExploreReps$.
\item The GitHub interface table of the Go benchmark suite~\cite{schmidt2020gohttproutingbench}, with
\GithubRoutes{} routes and four methods: one cell, $R=\ExGithubPairs$.
\item Long literals of \LongWordMin{} to \LongWordMax{} letters at $m=\SizeTenThousand$ and
\SizeHundredThousand, $R=\ExploreReps$. v2 keys a segment on its first \KeyBytes{} bytes, and short
literals might favour it.
\item A Zipf-distributed ring in the same cells, $R=\ExploreReps$.
\item \shape{mixed-overlap} without decoys at every size, $R=\ExNoDecoyPairs$, where gin takes part
because its walk then never needs to return.
\item The ablation arm in every cell of the main grid.
\end{itemize}
The pre-specification named one more run, for an H7 cell with a table seed apart from the others.
It measures that cell again, with the compile-time table copied into a heap block in the same
binary. Its condition held, and the run was not made in the frozen order. It was made afterwards,
as an exploratory run, gated by sanitizer records of its own (Section~\ref{sec:hseven}).

\subsection{Deviations and the Revision Log}\label{sec:deviation}

The revision log holds four entries. The first is a deviation, and it concerns the check on every
query at the end of the main grid. The frozen text fixed the check, but not how many checks run at
once. The run started
\VerifyJobsFirst{} at once. The actix-router checks at $m=\SizeHundredThousand$ took up to
\VerifyCheckGB~GB each and exhausted L's \VerifyMemoryGB~GB. The kernel killed \VerifyKilled{} of the
\VerifyCut{} checks, all after the last measured cell, and none had reported a wrong answer.

The repeat of all \VerifyCut{} checks, with \VerifyJobsRepeat{} at once, was logged before it ran,
because it could change H1's decision. It used the same checking binary from the run's own build,
after its hash was checked, and wrote a separate file. The first file was kept. The repeat checked
answers only: nothing was timed again, and no time was replaced. The later exploratory grids ran their
checks \VerifyJobsRepeat{} at once. The repeat changed which cells could be tested, as
Section~\ref{sec:hone} shows.

The other three entries change no decision. The second records that the paper's numbers come from two
generated files. One holds the analysis's macros; the other was written after the run, from the
analysis's outputs and the raw data. The third records that the conditional run of
Section~\ref{sec:explore} was due, was not made in the frozen order, and was made afterwards. The
fourth records how the
paper mentions the development round: the origin of H6's bound, and the direction of one cell,
without its numbers. One more departure is in the lab journal. A few seconds of analysis ran on L at
the lowest priority while the main grid measured. No process was excluded for it.
